\documentclass[12pt]{article}


% Codificación de caracteres e idioma
\usepackage[utf8]{inputenc}
\usepackage[spanish]{babel}

% Tamaño de página y márgenes
\usepackage[letterpaper, margin=3cm]{geometry}

% Orientación horizontal de página e hipervínculos
\usepackage{lscape, hyperref}

% Comillas, columnas, gráficos, símbolos matemáticos e ingreso de código de programación
\usepackage{csquotes, multicol, graphicx, amssymb, amsmath, listings, xcolor}

% Para tablas grandes y líneas estéticas de tablas
\usepackage{xltabular, booktabs}

% Para hacer mosaico de figuras
\usepackage{subcaption}

% Interlineado 1.5
\usepackage{setspace}
\onehalfspacing

% Uso de punto decimal, en lugar de coma
\decimalpoint

% Separación de columnas y sangría
\setlength{\columnsep}{2cm} 
\setlength{\parindent}{0cm}

% Páginado
\pagestyle{plain}

% Evitar separación de palabras
\hyphenpenalty=10000
\exhyphenpenalty=10000
\sloppy


\begin{document}

\begin{titlepage}
	\centering
	{\scshape\LARGE Universidad Aztlán \par}
	\vspace{1cm}
	{\scshape\Large Campus Playa del Carmen \par}
	\vspace{1.75cm}
	{\scshape\Huge Ecuaciones Diferenciales de Primer Orden \par}
    \vfill
    {\itshape\Large Teoría y Ejemplos \par}
	\vspace{2cm}
    {\itshape\Large Presentado por: \par}
	{\Large José Rosendo Camal Teh \par}
	\vfill
	{\Large Ingeniería en Sistemas Computacionales \par}
	{\Large Cuarto Cuatrimestre \par}
	{\Large Ecuaciones Diferenciales \par}
	\vfill
	{\Large Prof. Marcelino Galván Espinosa \par}
	\vfill
	{\Large 08 de noviembre de 2026 \par}
	\vfill
	{\Large Playa del Carmen, Quintana Roo \par}
\end{titlepage}
\newpage

\tableofcontents
\newpage 

\section*{\centering 10 EJERCICIOS DE NIVEL BÁSICO }
\addcontentsline{toc}{section}{10 EJERCICIOS DE NIVEL BÁSICO}

% Ejercicio 1
\subsection*{\centering Ejercicio 01 }
\addcontentsline{toc}{subsection}{Ejercicio 01}

\begin{itemize}
    \item \emph{Verifica si \(y = e^{2x}\) es una solución de la ecuación diferencial \(\frac{dy}{dx} = 2y\).}
\end{itemize}

Sea la ecuación
\[ y = e^{2x}. \]

Tenemos que
\begin{align*}
    \frac{d}{dx}\left( e^{2x} \right) &= e^{2x} \cdot \frac{d}{dx}\left( 2x \right) \\
    &= e^{2x} \cdot 2 \\
    &= 2e^{2x}.
\end{align*}

Dado que \(y = e^{2x}\), sustituimos:
\begin{align*}
    \frac{dy}{dx} &= 2y \\
    \frac{d}{dx}\left( e^{2x} \right) &= 2(e^{2x}).
\end{align*}

Por lo tanto, \(y = e^{2x}\) sí es una solución de \(\frac{dy}{dx} = 2y\).


% Ejercicio 2
\subsection*{\centering Ejercicio 02 }
\addcontentsline{toc}{subsection}{Ejercicio 02}

\begin{itemize}
    \item \emph{Determina el orden y el grado de la ecuación diferencial: \(\left(\frac{dy}{dx}\right)^{3} + 2y = x^{2}\).}
\end{itemize}

En esta ecuación, el término \(\left( \frac{dy}{dx} \right)^{3}\) contiene la única derivada de la expresión que está elevado al cubo y es una primera derivada. Por definición de orden y grado, la ecuación diferencial es de \textbf{primer orden} y de \textbf{tercer grado}.

% Ejercicio 3
\subsection*{\centering Ejercicio 03 }
\addcontentsline{toc}{subsection}{Ejercicio 03}

\begin{itemize}
    \item \emph{Determina el orden y el grado de la ecuación diferencial: \(\frac{d^{2}y}{dx^{2}} + 3\frac{dy}{dx} = \sen{x}\).}
\end{itemize}

En esta ecuación diferencial, el término \(\frac{d^{2}y}{dx^{2}}\) contiene la derivada más alta (con primera potencia) y es una segunda derivada. Por lo tanto, la ecuación presente es de \textbf{segundo orden} y de \textbf{primer grado}.

% Ejercicio 4
\subsection*{\centering Ejercicio 04 }
\addcontentsline{toc}{subsection}{Ejercicio 04}

\begin{itemize}
    \item \emph{¿La ecuación diferencial \(\frac{dy}{dx} + y^{2} = x\) es lineal o no lineal? Justifica tu respuesta.}
\end{itemize}

Es una \textbf{ecuación no lineal} porque a pesar que el término \(\frac{dy}{dx}\) está elevado a la primera potencia y no tiene como uno de sus factores la función \(y\), el término \(y^{2}\) no está elevado a la potencia 1 y esto incumple la definición de ecuación lineal.

% Ejercicio 5
\subsection*{\centering Ejercicio 05 }
\addcontentsline{toc}{subsection}{Ejercicio 05}

\begin{itemize}
    \item \emph{Encuentra la solución general de la ecuación diferencial \(\frac{dy}{dx} = 3x^{2}\).}
\end{itemize}

Hallar la solución general de esta ecuación implica buscar una función tal que su derivada sea \(3x^{2}\). En otros términos, la solución es la integral de \(3x^{2}\):
\begin{align*}
    \int 3x^{2} \, dx &= 3 \left( \frac{x^{2 + 1}}{2 + 1} \right) + C \\
    &= 3 \left( \frac{x^{3}}{3} \right) + C \\
    &= x^{3} + C
\end{align*}

Dado que la integral y la derivada son operaciones inversas entre sí, se concluye que la solución general es:
\[ y = x^{3} + C. \]

% Ejercicio 6
\subsection*{\centering Ejercicio 06 }
\addcontentsline{toc}{subsection}{Ejercicio 06}

\begin{itemize}
    \item \emph{Encuentra la solución general de la ecuación diferencial \(\frac{dy}{dx} = 5y\).}
\end{itemize}

Hallar la solución general implica encontrar una función tal que su derivada sea \(5y\). Un planteamiento similar es \(\frac{dy}{dx} = y\) cuya solución es \(y = Ce^{x}\) y como la derivada de \(e^{ax}\) es \(ae^{ax}\) derivamos \(y = e^{5x}\):
\begin{align*}
    \frac{dy}{dx} &= \frac{d}{dx}\left( e^{5x} \right) \\
    &= e^{5x} \cdot \frac{d}{dx}\left( 5x \right) \\
    &= 5e^{5x}
\end{align*}

Por lo que se confirma la solución general \(y = Ce^{5x}\), siendo \(y = 5e^{5x}\) una solución particular.

% Ejercicio 7
\subsection*{\centering Ejercicio 07 }
\addcontentsline{toc}{subsection}{Ejercicio 07}

\begin{itemize}
    \item \emph{Resuelve el siguiente Problema de Valor Inicial: \(\frac{dy}{dx} = 4x\), con la condición inicial \(y(1) = 3\).}
\end{itemize}

Primero, se obtiene la solución general:
\begin{align*}
    \int 4x \, dx &= 4\frac{x^{2}}{2} + C \\
    &= 2x^{2} + C \\
    y &= 2x^{2} + C.
\end{align*}

Después, se sustituye \(x = 1\) y \(y = 3\) en \(y = 2x^{2} + C\):
\begin{align*}
    3 &= 2(1)^{2} + C \\
    3 &= 2 + C \\
    2 + C &= 3 \\
    C &= 3 - 2 \\
    C &= 1
\end{align*}

Por lo tanto, la solución particular es \(y = 2x^{2} + 1\).

% Ejercicio 8
\subsection*{\centering Ejercicio 08 }
\addcontentsline{toc}{subsection}{Ejercicio 08}

\begin{itemize}
    \item \emph{Resuelve el siguiente Problema de Valor Inicial: \(\frac{dy}{dx} = -2y\), con la condición inicial \(y(0) = 4\).}
\end{itemize}

Por similitud a los problemas anteriores, se determina la solución general
\[ y = Ce^{-2x}. \]

Se sustituye \(x = 0\) y \(y = 4\) en \(y = Ce^{-2x}\):
\begin{align*}
    4 &= Ce^{-2 \cdot 0} \\
    4 &= Ce^{0} \\
    4 &= C \cdot 1 \\
    C &= 4
\end{align*}

Por lo tanto, la solución particular es \(y = 4e^{-2x}\).

% Ejercicio 9
\subsection*{\centering Ejercicio 09 }
\addcontentsline{toc}{subsection}{Ejercicio 09}

\begin{itemize}
    \item \emph{Verifica si el PVI \(\frac{dy}{dx} = x + y\), \(y(0) = 2\) cumple con las condiciones del Teorema de Existencia y Unicidad en un rectángulo alrededor del punto \((0, 2)\).}
\end{itemize}

Sea la función \(f(x, y) = x + y\) y su derivada parcial \(\frac{\partial f}{\partial y} = 1\) tenemos que la función \(f\) es continua en todo el plano al igual que su derivada parcial. Por lo anterior, se concluye que hay solución cerca de \(x = 0\).

% Ejercicio 10
\subsection*{\centering Ejercicio 10 }
\addcontentsline{toc}{subsection}{Ejercicio 10}

\begin{itemize}
    \item \emph{Verifica si el PVI \(\frac{dy}{dx} = \sqrt{y}\), \(y(0) = 0\) cumple con las condiciones del Teorema de Existencia y Unicidad. ¿Qué concluyes?}
\end{itemize}

Sea la función \(f(x, y) = \sqrt{y}\) y su derivada parcial:
\begin{align*}
    \frac{\partial f}{\partial y} &= \frac{\partial}{\partial y}\left(\sqrt{y}\right) \\
    &= \frac{\partial}{\partial y}\left(y^{\frac{1}{2}}\right) \\
    &= \frac{1}{2\sqrt{y}}.
\end{align*}

Dado que la derivada parcial no es continua cuando \(y = 0\) el teorema no garantiza existencia ni unicidad.

\end{document}